\documentclass[11pt]{article}
\usepackage[utf8]{inputenc}
\usepackage{amsmath, amssymb}
\usepackage{geometry}
\usepackage{listings}
\usepackage{xcolor}
\usepackage{tabularx}
\usepackage{booktabs}
\usepackage{hyperref}

\geometry{a4paper, margin=1in}

\definecolor{codegray}{rgb}{0.5,0.5,0.5}
\definecolor{codepurple}{rgb}{0.58,0,0.82}
\definecolor{backcolour}{rgb}{0.95,0.95,0.92}

\lstset{
    language=Python,
    backgroundcolor=\color{backcolour},
    commentstyle=\color{codegray},
    keywordstyle=\color{magenta},
    stringstyle=\color{codepurple},
    basicstyle=\ttfamily\small,
    breaklines=true,
    showstringspaces=false
}

\title{Abstract Algebra Lab: Cayley Graphs \& Structural Group Theory \\ \large Problem Set \& Implementation Guide}
\author{VB}
\date{\today}

\begin{document}

\maketitle

\section*{Problem 0: Start the engine}

Consider the files given.
\begin{center}
\renewcommand{\arraystretch}{1.5}
\begin{tabularx}{\textwidth}{@{} >{\bfseries\raggedright\arraybackslash}p{5.5cm} X @{}}
\toprule
\textcolor{blue}{Filename} & \textcolor{blue}{Purpose} \\ \midrule
\multicolumn{2}{@{}l}{\textbf{Direct Interaction}} \\ \midrule
\lstinline|implementation_tasks.py| & \textcolor{red}{Here you write your code.} Contains the core group theory algorithms (Cayley Graph BFS, Subgroup Generation, Cosets, Conjugacy) intended for student implementation. You can add your own \lstinline|unittest.TestCase| classes at the bottom and run \lstinline|python implementation_tasks.py| to execute them. You can add your own \lstinline|unittest.TestCase| classes at the bottom and run \lstinline|python implementation_tasks.py| to execute them. \\
\lstinline|lab_dashboard.py| & \textcolor{red}{This file you run.} A dedicated educational sandbox for debugging individual algebra levels through an interactive UI. \\
\lstinline|levels/level_*.py| & Graphical representation for each specific problem. \textbf{Can be run standalone.} You can use \lstinline|python levels/level_*.py| for the graphical debugger. \\ \midrule
\multicolumn{2}{@{}l}{\textbf{Under the Hood}} \\ \midrule
\lstinline|group_engine.py| & Wrapper for the SymPy Todd-Coxeter Finitely Presented Group evaluator. Provides a unified \lstinline|GroupElement| object with overloaded operators (\lstinline|*|, \lstinline|inv()|, \lstinline|==|) to bypass manual string reduction and solve the Word Problem mathematically.\\
\lstinline|levels/base_level.py| & Base class for graphical debuggers, managing matplotlib node layouts and interaction.\\
\bottomrule
\end{tabularx}
\end{center}

Run the file \lstinline|lab_dashboard.py|. You will see a window with a menu of group theory problems.

Now use any editor to open the file \lstinline|implementation_tasks.py|. You should see the API for functions you are expected to implement.
\begin{lstlisting}
import numpy as np

# =====================================================================
# STUDENT IMPLEMENTATION (Group Theory & Cayley Graphs)
# =====================================================================

def generate_cayley_graph(group, generators):
    pass

def generate_subgroup(group, subset_strings):
    pass

# ... (MORE CODE)
\end{lstlisting}

\subsubsection*{Self-Testing (Optional)}
At the bottom of \lstinline|implementation_tasks.py|, you will find a block that looks like this:
\begin{lstlisting}[language=Python]
if __name__ == '__main__':
    import unittest
    # Add your own unittest.TestCase classes here...
    unittest.main(verbosity=2)
\end{lstlisting}
This allows you to write your own automated unit tests to verify your code before running the visual dashboard. If you'd like to test your functions with specific inputs, simply define a class inheriting from \lstinline|unittest.TestCase| above the \lstinline|unittest.main()| call. For example, if you wanted to test simple math:
\begin{lstlisting}[language=Python]
class TestMyMath(unittest.TestCase):
    def test_addition(self):
        result = 2 + 2
        self.assertEqual(result, 4, "Addition is broken!")
\end{lstlisting}
You can run these tests anytime by executing \lstinline|python implementation_tasks.py| in your terminal. This is completely optional, but highly recommended for debugging tricky edge cases.

Now proceed with the following problems. The \lstinline|GroupElement| objects you are passed support native mathematical operations. You can multiply two elements with \lstinline|a * b|, find an inverse with \lstinline|a.inv()|, and test equivalence with \lstinline|a == b|. You do not need to perform string reduction yourself.

\section*{Problem 1: Cayley Graph Generation}

\subsubsection*{Mathematical part}
\textbf{Given:} A finite group $G$ and a set of generator symbols $S$. \\
\textbf{Derive:} The Cayley Graph $\Gamma(G, S)$, a directed graph where nodes correspond to elements of $G$, and directed edges $(g, s, g \cdot s)$ exist for every $g \in G, s \in S$. The Cayley graph visualizes the connectivity of the group structure.

\subsubsection*{Implementation part}
\textbf{Implement:} \lstinline|generate_cayley_graph(group, generators)|. Use a Breadth-First Search (BFS) starting from the identity $e$. For each popped element, iterate over all generators, multiply them on the right, and add newly discovered nodes to the queue.
\begin{itemize}
    \item \textbf{Returns:} \lstinline|(nodes, edges)| where \lstinline|nodes| is a \lstinline|set| of unique \lstinline|GroupElements|, and \lstinline|edges| is a list of tuples: \lstinline|(source, generator_string, target)|.
\end{itemize}

\section*{Problem 2: Subgroups and Cosets}

\subsubsection*{Mathematical part}
\textbf{Given:} A subset of elements $H \subseteq G$. \\
\textbf{Derive:} The full subgroup $\langle H \rangle$ generated by $H$. This requires closure under multiplication and inversion. Then, derive all left cosets of $\langle H \rangle$ in $G$. A left coset is defined as $gH = \{g \cdot h \mid h \in \langle H \rangle \}$.

\subsubsection*{Implementation part}
\textbf{Implement:}
\begin{itemize}
    \item \lstinline|generate_subgroup(group, subset_strings)|: Use a BFS queue to generate closure under multiplication and inversion for the given subset of generator strings. Return a set of elements.
    \item \lstinline|compute_left_cosets(group, subgroup_elements)|: Given a fully generated subgroup, partition the group into left cosets. Note that $g_1 H = g_2 H$ if $g_2^{-1} g_1 \in H$. Return a list of sets. Ensure the subgroup itself is the first coset in the list.
\end{itemize}

\section*{Problem 3: Normality and Conjugacy}

\subsubsection*{Mathematical part}
\textbf{Given:} A group $G$ and an element $x \in G$. \\
\textbf{Derive:} The conjugacy class of $x$, defined as $Cl(x) = \{ g \cdot x \cdot g^{-1} \mid g \in G \}$. \\
\textbf{Given:} A subgroup $H \le G$. \\
\textbf{Derive:} Whether $H$ is normal in $G$ ($H \triangleleft G$). A subgroup is normal if it is invariant under conjugation: $g \cdot h \cdot g^{-1} \in H$ for all $g \in G, h \in H$.

\subsubsection*{Implementation part}
\textbf{Implement:}
\begin{itemize}
    \item \lstinline|compute_conjugacy_class(group, element)|: Return a set of elements representing the conjugacy class.
    \item \lstinline|compute_normal_subgroups(group, list_of_subgroups)|: Given a list of subgroup sets, return a list of booleans indicating whether each subgroup is strictly normal.
\end{itemize}

\section*{Problem 4: Center and Commutator}

\subsubsection*{Mathematical part}
\textbf{Given:} A group $G$. \\
\textbf{Derive:} 
\begin{itemize}
    \item The Center $Z(G) = \{ z \in G \mid z \cdot g = g \cdot z \text{ for all } g \in G\}$.
    \item The Commutator Subgroup $[G, G]$, the group generated by all commutators $[x, y] = x \cdot y \cdot x^{-1} \cdot y^{-1}$ for $x,y \in G$.
\end{itemize}

\subsubsection*{Implementation part}
\textbf{Implement:}
\begin{itemize}
    \item \lstinline|compute_center(group)|: Iterate over all elements and find those that commute with every other element. Return as a set.
    \item \lstinline|compute_commutator_subgroup(group)|: First compute all simple commutators. Then, since the set of simple commutators is not always closed, use BFS to find the full generated subgroup. Return as a set.
\end{itemize}

\end{document}
